\documentclass[11pt]{article}
\usepackage{docstyle}
\renewcommand\sheetname{Homework 1 \textperiodcentered{} Green Machines}
\begin{document}
\sheethead{Homework 1: Green Machines}{Year 9 Science \textperiodcentered{} about 40 minutes}{Section A: A \quad Section B: M / E}

\section{Section A}
\Q Write the meaning of each word.\lvl{A}
\par\sub{a} germination
\alines{1}{when a root and shoot grow out of a seed and it develops into a seedling}
\sub{b} transpiration
\alines{1}{the loss of water vapour from a plant's leaves through the stomata}
\sub{c} zygote
\alines{1}{the cell formed when the pollen nucleus fuses with the ovule nucleus}

\Q Complete the sentences about flower parts.\lvl{A}
\par The \blank[26mm]{sepal} protects the flower while it is a bud. The \blank[26mm]{nectary} makes sweet nectar.
The \blank[26mm]{stigma} is sticky so that pollen grains stay on it.\par\vspace{2mm}

\Q Name three things a seedling competes with its parent plant for.\lvl{A}
\par 1. \blank[34mm]{light} \quad 2. \blank[34mm]{water} \quad 3. \blank[34mm]{minerals (or space)}\par\vspace{2mm}

\section{Quick review}
\Q On a hot, dry afternoon a gardener sees the leaves of her tomato plants drooping.\lvl{A}
\par\sub{a} Which process is the plant losing water by? \hfill\blank[40mm]{transpiration}
\par\sub{b} Through which openings in the leaf is the water lost? \hfill\blank[40mm]{stomata}
\par\sub{c} Which tubes should be bringing water up from the roots? \hfill\blank[40mm]{xylem}\par\vspace{1mm}

\section{Section B}
\Q Explain how a bee pollinates a flower. Use the names of the flower parts.\lvl{M}
\alines{4}{The bee is attracted by the brightly coloured petals, the scent and the nectar. Sticky pollen from the anthers rubs onto its body. When it visits another flower of the same species, the pollen rubs off onto the stigma. This is cross-pollination.}

\Q Explain why cross-pollination is better for a species than self-pollination.\lvl{E}
\alines{4}{Cross-pollination mixes the genetic material of two parents, so the offspring show more variation. If the environment changes or a disease arrives, some offspring are likely to have features that let them survive and reproduce, so the species continues.}

\Q Gardeners often soak bean seeds in water overnight before planting them. Explain why.\lvl{M}
\alines{3}{Water softens the testa, so the radicle can grow through it. Water is also one of the conditions for germination, so the seeds start to germinate sooner.}

\Q A carrot left in fresh water for a day becomes firm. A carrot left in very salty water becomes soft and bendy. Explain both results using osmosis.\lvl{E}
\alines{5}{Fresh water has a higher concentration of water than the carrot cells, so water moves into the cells by osmosis through the semi-permeable membranes and the carrot becomes firm. The salty water has a lower concentration of water than the cells, so water moves out of the cells by osmosis and the carrot goes soft.}

\Q A potted plant is weighed each day. All the mass it loses is water.\lvl{M}
\par\vspace{1mm}
\begin{tabular}{@{}l c c c c@{}}
\toprule
Day & 1 & 2 & 3 & 4\\
Mass (g) & 500 & 488 & 476 & 464\\
\bottomrule
\end{tabular}\par\vspace{1mm}
\sub{a} How much water does the plant lose each day?
\alines{1}{12 g per day ($500 - 488 = 12$ g)}
\sub{b} Explain where the water goes.
\alines{2}{It evaporates from the leaves and is lost as water vapour through the stomata (transpiration).}

\Q Sam explains how water gets to the leaves of a plant.\lvl{E}
\par\vspace{1mm}
\samcard{0.9\linewidth}{\flawed{Water enters the root hairs by osmosis, from a low concentration of water to a high concentration of water. It travels up the phloem to the leaves. The leaves make glucose and oxygen from carbon dioxide and water. At night the plant stops respiring.}}\par\vspace{1mm}
\textbf{Sam's answer has 3 errors.} Find each one, fix it, and say why it loses marks.
\alines{5}{1. Osmosis is from a high to a low concentration of water (\Lref{osmosis}). 2. Water travels up the xylem, not the phloem (\Lref{xylem}). 3. Plants respire all the time, day and night (\Lref{respire}).}

\vspace{4mm}
{\small Also set: Green Machines booklet p.24, fill in the glossary for any ten words.}
\end{document}
